\documentclass[11pt]{article}
\usepackage[utf8]{inputenc}
\usepackage{amsmath, amssymb, amsthm}
\usepackage{geometry}
\geometry{letterpaper, margin=1in}
\usepackage{hyperref}
\usepackage{enumitem}

\title{\textbf{Homework 2.5: Differential Equations \\ Exact Equations \& Integrating Factors} \\ \Large Problem Sets \& Solutions}
\author{}
\date{}

\begin{document}

\maketitle
\tableofcontents
\newpage

% ---------------------------------------------------------
\section{Concept Checklist \& Question Mapping}
% ---------------------------------------------------------
Based on the provided source material, this section outlines the core concepts, tricks, and required skills to master Exact Equations and Integrating Factors. 

\begin{itemize}
    \item \textbf{C1: Verification of Exactness} \\
    Identifying standard form $M dx + N dy = 0$ and checking if $M_y = N_x$. \\
    \textit{Mapped Questions:} 1, 2, 3, 4
    
    \item \textbf{C2: Solving Standard Exact Equations} \\
    Constructing the implicit solution $F(x,y) = C$ by integrating $M$ with respect to $x$ and $N$ with respect to $y$, taking the union of unique terms. \\
    \textit{Mapped Questions:} 5, 6, 7, 8
    
    \item \textbf{C3: Finding 1D Integrating Factors ($\mu(x)$ or $\mu(y)$)} \\
    Utilizing the formulas $\frac{M_y - N_x}{N} = h(x)$ or $\frac{N_x - M_y}{M} = k(y)$ to find the proper integrating factor. \\
    \textit{Mapped Questions:} 9, 10, 11, 12
    
    \item \textbf{C4: Solving Non-Exact Equations via 1D Integrating Factor} \\
    Applying the derived 1D integrating factor to force exactness, then solving. \\
    \textit{Mapped Questions:} 13, 14, 15, 16, 17, 18
    
    \item \textbf{C5: Algebraic Integrating Factors of form $\mu = x^n y^m$} \\
    Multiplying literal terms through the equation and setting $M'_y = N'_x$ to form a system of linear equations to solve for $n$ and $m$. \\
    \textit{Mapped Questions:} 19, 20, 21, 22
    
    \item \textbf{C6: Standard First-Order Linear ODEs} \\
    Rearranging into $y' + P(x)y = Q(x)$ and using the factor $e^{\int P(x) dx}$. \\
    \textit{Mapped Questions:} 23, 24, 25
    
    \item \textbf{C7: Bernoulli Differential Equations} \\
    Applying the substitution $v = y^{1-n}$ to transform a non-linear Bernoulli equation into a standard linear ODE. \\
    \textit{Mapped Questions:} 26, 27, 28
    
    \item \textbf{C8: Identifying Lost or Gained Solutions} \\
    Recognizing that multiplying by integrating factors that restrict the domain (e.g., $y^{-2}$) can cause valid trivial solutions to be "lost." \\
    \textit{Mapped Questions:} 29
    
    \item \textbf{C9: Diagnostic Testing (``Find the Flaw'')} \\
    Identifying common student pitfalls, such as failing to establish standard form before finding $P(x)$, or adding overlapping integral terms. \\
    \textit{Mapped Questions:} 30
\end{itemize}

\newpage
% ---------------------------------------------------------
\section{Problem Sets}
% ---------------------------------------------------------

\subsection*{C1: Verification of Exactness}
\textit{Determine whether the following differential equations are exact. Justify your answer.}

\begin{enumerate}
    \item $(2xy^2 + 3x^2) dx + (2x^2y + 4y^3) dy = 0$
    \item $(e^x \sin y + 2x) dx + (e^x \cos y + y^2) dy = 0$
    \item $(3x^2 y - y^3) dx + (x^3 - 3xy^2 + y) dy = 0$
    \item $(\ln y) dx + \left(\frac{x}{y}\right) dy = 0$
\end{enumerate}

\subsection*{C2: Solving Standard Exact Equations}
\textit{Verify that the equations are exact and find the implicit solution $F(x,y) = C$.}

\begin{enumerate}[resume]
    \item $(2xy + 3) dx + (x^2 - 1) dy = 0$
    \item $(\cos x \cos y + 2x) dx - (\sin x \sin y + 2y) dy = 0$
    \item $\left(y e^{xy} - \frac{1}{x}\right) dx + \left(x e^{xy} + \frac{1}{y}\right) dy = 0$
    \item $(3x^2 y^2 + 2x) dx + (2x^3 y - 4y^3) dy = 0$
\end{enumerate}

\subsection*{C3: Finding 1D Integrating Factors}
\textit{The following equations are not exact. Find an integrating factor that is a function of only $x$ or only $y$.}

\begin{enumerate}[resume]
    \item $(3xy + 2y^2) dx + (x^2 + 2xy) dy = 0$
    \item $(4xy) dx + (3x^2 + y^2) dy = 0$
    \item $(x^2 + y^2 + x) dx + (xy) dy = 0$
    \item $(3y) dx + (2x) dy = 0$
\end{enumerate}

\subsection*{C4: Solving Non-Exact Equations via 1D Integrating Factor}
\textit{Find an appropriate integrating factor, multiply the ODE by it, and solve the newly exact equation.}

\begin{enumerate}[resume]
    \item $(3xy + y^2) dx + (x^2 + xy) dy = 0$
    \item $(2y^2 + 3x) dx + (2xy) dy = 0$
    \item $(y) dx + (y^2 - x) dy = 0$
    \item $(2xy) dx + (y^2 - 3x^2) dy = 0$
    \item $(x + y^2) dx - (2xy) dy = 0$
    \item $(2xy^2 + y) dx + (2y^3 - x) dy = 0$
\end{enumerate}

\subsection*{C5: Algebraic Integrating Factors of Form $\mu = x^n y^m$}
\textit{Find an integrating factor of the form $x^n y^m$ and solve the equation.}

\begin{enumerate}[resume]
    \item $(3xy + 2y^3) dx + (2x^2 + 4xy^2) dy = 0$
    \item $(2xy - y^2) dx + (3x^2 - 4xy) dy = 0$
    \item $(2x^3 y - y^4) dx + (2xy^3 - x^4) dy = 0$
    \item $(2y^2 + 3xy) dx + (4xy + 3x^2) dy = 0$
\end{enumerate}

\subsection*{C6: Standard First-Order Linear ODEs}
\textit{Solve using the standard linear integrating factor $\mu(x) = e^{\int P(x)dx}$.}

\begin{enumerate}[resume]
    \item $\frac{dy}{dx} - 3y = e^{3x}$
    \item $x \frac{dy}{dx} + 2y = \sin(x)$ \quad (for $x > 0$)
    \item $\frac{dy}{dx} + y \cot x = \sec x$
\end{enumerate}

\subsection*{C7: Bernoulli Differential Equations}
\textit{Use the substitution $v = y^{1-n}$ to reduce to a linear ODE, then solve.}

\begin{enumerate}[resume]
    \item $\frac{dy}{dx} + \frac{1}{x}y = x y^2$
    \item $x \frac{dy}{dx} + y = y^{-2}$
    \item $\frac{dy}{dx} - y = e^x y^2$
\end{enumerate}

\subsection*{C8: Identifying Lost or Gained Solutions}
\begin{enumerate}[resume]
    \item Find an integrating factor $\mu(y)$ to solve the equation $(y^3) dx + (xy^2) dy = 0$. After solving, analyze your steps and clearly state whether any trivial solutions to the original ODE were lost during the process.
\end{enumerate}

\subsection*{C9: Diagnostic Testing (``Find the Flaw'')}
\begin{enumerate}[resume]
    \item \textbf{Problem:} Solve $x \frac{dy}{dx} + 2y = x^2$ using an integrating factor. \\
    \textbf{Flawed Solution:} The equation is of the form $y' + P(x)y = Q(x)$ where $P(x) = 2$. The integrating factor is $\mu(x) = e^{\int 2 dx} = e^{2x}$. Multiplying by $e^{2x}$ yields $x e^{2x} y' + 2 e^{2x} y = x^2 e^{2x}$. Integrating both sides will give the solution. \\
    \textbf{Task:} Act as the grader. In one sentence, identify the mathematical flaw in the student's reasoning, then provide the correct initial step to solve the problem.
\end{enumerate}

\newpage
% ---------------------------------------------------------
\section{Detailed Solutions}
% ---------------------------------------------------------

\subsection*{C1: Verification of Exactness}

\textbf{1.} $M = 2xy^2 + 3x^2$, $N = 2x^2y + 4y^3$ \\
$M_y = \frac{\partial}{\partial y}(2xy^2 + 3x^2) = 4xy$ \\
$N_x = \frac{\partial}{\partial x}(2x^2y + 4y^3) = 4xy$ \\
Since $M_y = N_x$, the equation is \textbf{exact}.

\vspace{0.5cm}
\textbf{2.} $M = e^x \sin y + 2x$, $N = e^x \cos y + y^2$ \\
$M_y = e^x \cos y$ \\
$N_x = e^x \cos y$ \\
Since $M_y = N_x$, the equation is \textbf{exact}.

\vspace{0.5cm}
\textbf{3.} $M = 3x^2 y - y^3$, $N = x^3 - 3xy^2 + y$ \\
$M_y = 3x^2 - 3y^2$ \\
$N_x = 3x^2 - 3y^2$ \\
Since $M_y = N_x$, the equation is \textbf{exact}.

\vspace{0.5cm}
\textbf{4.} $M = \ln y$, $N = \frac{x}{y}$ \\
$M_y = \frac{1}{y}$ \\
$N_x = \frac{1}{y}$ \\
Since $M_y = N_x$, the equation is \textbf{exact}.

\subsection*{C2: Solving Standard Exact Equations}

\textbf{5.} $(2xy + 3) dx + (x^2 - 1) dy = 0$ \\
$M_y = 2x$, $N_x = 2x \implies$ Exact. \\
$\int M dx = \int (2xy + 3) dx = x^2 y + 3x$ \\
$\int N dy = \int (x^2 - 1) dy = x^2 y - y$ \\
Union of unique terms: $F(x,y) = x^2 y + 3x - y$. \\
\textbf{Solution:} $x^2 y + 3x - y = C$.

\vspace{0.5cm}
\textbf{6.} $(\cos x \cos y + 2x) dx - (\sin x \sin y + 2y) dy = 0$ \\
$M_y = -\cos x \sin y$, $N_x = -\cos x \sin y \implies$ Exact. \\
$\int M dx = \sin x \cos y + x^2$ \\
$\int N dy = \int (-\sin x \sin y - 2y) dy = \sin x \cos y - y^2$ \\
Union of unique terms: $F(x,y) = \sin x \cos y + x^2 - y^2$. \\
\textbf{Solution:} $\sin x \cos y + x^2 - y^2 = C$.

\vspace{0.5cm}
\textbf{7.} $\left(y e^{xy} - \frac{1}{x}\right) dx + \left(x e^{xy} + \frac{1}{y}\right) dy = 0$ \\
$M_y = e^{xy} + xye^{xy}$, $N_x = e^{xy} + xye^{xy} \implies$ Exact. \\
$\int M dx = e^{xy} - \ln|x|$ \\
$\int N dy = e^{xy} + \ln|y|$ \\
Union of unique terms: $F(x,y) = e^{xy} - \ln|x| + \ln|y|$. \\
\textbf{Solution:} $e^{xy} - \ln|x| + \ln|y| = C$.

\vspace{0.5cm}
\textbf{8.} $(3x^2 y^2 + 2x) dx + (2x^3 y - 4y^3) dy = 0$ \\
$M_y = 6x^2 y$, $N_x = 6x^2 y \implies$ Exact. \\
$\int M dx = x^3 y^2 + x^2$ \\
$\int N dy = x^3 y^2 - y^4$ \\
Union of unique terms: $F(x,y) = x^3 y^2 + x^2 - y^4$. \\
\textbf{Solution:} $x^3 y^2 + x^2 - y^4 = C$.

\subsection*{C3: Finding 1D Integrating Factors}

\textbf{9.} $(3xy + 2y^2) dx + (x^2 + 2xy) dy = 0$ \\
$M_y = 3x + 4y$, $N_x = 2x + 2y$. \\
$\frac{M_y - N_x}{N} = \frac{(3x+4y) - (2x+2y)}{x^2+2xy} = \frac{x+2y}{x(x+2y)} = \frac{1}{x}$. \\
Since it depends only on $x$, $\mu(x) = e^{\int \frac{1}{x} dx} = e^{\ln|x|} = x$. \\
\textbf{Integrating Factor:} $\mu(x) = x$.

\vspace{0.5cm}
\textbf{10.} $(4xy) dx + (3x^2 + y^2) dy = 0$ \\
$M_y = 4x$, $N_x = 6x$. \\
Check $\frac{N_x - M_y}{M} = \frac{6x - 4x}{4xy} = \frac{2x}{4xy} = \frac{1}{2y}$. \\
Since it depends only on $y$, $\mu(y) = e^{\int \frac{1}{2y} dy} = e^{\frac{1}{2}\ln y} = y^{1/2} = \sqrt{y}$. \\
\textbf{Integrating Factor:} $\mu(y) = \sqrt{y}$.

\vspace{0.5cm}
\textbf{11.} $(x^2 + y^2 + x) dx + (xy) dy = 0$ \\
$M_y = 2y$, $N_x = y$. \\
$\frac{M_y - N_x}{N} = \frac{2y - y}{xy} = \frac{y}{xy} = \frac{1}{x}$. \\
$\mu(x) = e^{\int \frac{1}{x} dx} = e^{\ln|x|} = x$. \\
\textbf{Integrating Factor:} $\mu(x) = x$.

\vspace{0.5cm}
\textbf{12.} $(3y) dx + (2x) dy = 0$ \\
$M_y = 3$, $N_x = 2$. \\
$\frac{M_y - N_x}{N} = \frac{3 - 2}{2x} = \frac{1}{2x}$. \\
$\mu(x) = e^{\int \frac{1}{2x} dx} = e^{\frac{1}{2}\ln x} = x^{1/2} = \sqrt{x}$. \\
\textbf{Integrating Factor:} $\mu(x) = \sqrt{x}$.

\subsection*{C4: Solving Non-Exact Equations via 1D Integrating Factor}

\textbf{13.} $(3xy + y^2) dx + (x^2 + xy) dy = 0$ \\
$M_y = 3x+2y, N_x = 2x+y$. $\frac{M_y - N_x}{N} = \frac{x+y}{x(x+y)} = \frac{1}{x} \implies \mu(x) = x$. \\
Multiply by $x$: $(3x^2y + xy^2) dx + (x^3 + x^2y) dy = 0$. \\
$\int M' dx = x^3 y + \frac{1}{2}x^2 y^2$. $\int N' dy = x^3 y + \frac{1}{2}x^2 y^2$. \\
\textbf{Solution:} $x^3 y + \frac{1}{2}x^2 y^2 = C$.

\vspace{0.5cm}
\textbf{14.} $(2y^2 + 3x) dx + (2xy) dy = 0$ \\
$M_y = 4y, N_x = 2y$. $\frac{M_y - N_x}{N} = \frac{2y}{2xy} = \frac{1}{x} \implies \mu(x) = x$. \\
Multiply by $x$: $(2xy^2 + 3x^2) dx + (2x^2y) dy = 0$. \\
$\int M' dx = x^2 y^2 + x^3$. $\int N' dy = x^2 y^2$. \\
\textbf{Solution:} $x^2 y^2 + x^3 = C$.

\vspace{0.5cm}
\textbf{15.} $(y) dx + (y^2 - x) dy = 0$ \\
$M_y = 1, N_x = -1$. $\frac{N_x - M_y}{M} = \frac{-2}{y} \implies \mu(y) = e^{\int \frac{-2}{y} dy} = y^{-2}$. \\
Multiply by $y^{-2}$: $(y^{-1}) dx + (1 - xy^{-2}) dy = 0$. \\
$\int M' dx = x y^{-1}$. $\int N' dy = y + x y^{-1}$. \\
\textbf{Solution:} $x y^{-1} + y = C$.

\vspace{0.5cm}
\textbf{16.} $(2xy) dx + (y^2 - 3x^2) dy = 0$ \\
$M_y = 2x, N_x = -6x$. $\frac{N_x - M_y}{M} = \frac{-8x}{2xy} = \frac{-4}{y} \implies \mu(y) = y^{-4}$. \\
Multiply by $y^{-4}$: $(2x y^{-3}) dx + (y^{-2} - 3x^2 y^{-4}) dy = 0$. \\
$\int M' dx = x^2 y^{-3}$. $\int N' dy = -y^{-1} + x^2 y^{-3}$. \\
\textbf{Solution:} $x^2 y^{-3} - y^{-1} = C$.

\vspace{0.5cm}
\textbf{17.} $(x + y^2) dx - (2xy) dy = 0$ \\
$M_y = 2y, N_x = -2y$. $\frac{M_y - N_x}{N} = \frac{4y}{-2xy} = \frac{-2}{x} \implies \mu(x) = x^{-2}$. \\
Multiply by $x^{-2}$: $(x^{-1} + x^{-2}y^2) dx - (2x^{-1}y) dy = 0$. \\
$\int M' dx = \ln|x| - x^{-1}y^2$. $\int N' dy = -x^{-1}y^2$. \\
\textbf{Solution:} $\ln|x| - x^{-1}y^2 = C$.

\vspace{0.5cm}
\textbf{18.} $(2xy^2 + y) dx + (2y^3 - x) dy = 0$ \\
$M_y = 4xy+1, N_x = -1$. $\frac{N_x - M_y}{M} = \frac{-4xy - 2}{y(2xy+1)} = \frac{-2(2xy+1)}{y(2xy+1)} = \frac{-2}{y} \implies \mu(y) = y^{-2}$. \\
Multiply by $y^{-2}$: $(2x + y^{-1}) dx + (2y - xy^{-2}) dy = 0$. \\
$\int M' dx = x^2 + xy^{-1}$. $\int N' dy = y^2 + xy^{-1}$. \\
\textbf{Solution:} $x^2 + xy^{-1} + y^2 = C$.

\subsection*{C5: Algebraic Integrating Factors of Form $\mu = x^n y^m$}

\textbf{19.} $(3xy + 2y^3) dx + (2x^2 + 4xy^2) dy = 0$ \\
Multiply by $x^n y^m$: \\
$M' = 3x^{n+1}y^{m+1} + 2x^n y^{m+3} \implies M'_y = 3(m+1)x^{n+1}y^m + 2(m+3)x^n y^{m+2}$. \\
$N' = 2x^{n+2}y^m + 4x^{n+1}y^{m+2} \implies N'_x = 2(n+2)x^{n+1}y^m + 4(n+1)x^n y^{m+2}$. \\
Equate coefficients: \\
$3m+3 = 2n+4 \implies 2n - 3m = -1$ \\
$2m+6 = 4n+4 \implies 4n - 2m = 2 \implies 2n - m = 1$ \\
Subtracting the first from the second: $2m = 2 \implies m = 1$. Then $2n - 1 = 1 \implies n = 1$. \\
Integrating factor is $\mu = xy$. Multiply original by $xy$: \\
$(3x^2 y^2 + 2xy^4) dx + (2x^3 y + 4x^2 y^3) dy = 0$. \\
\textbf{Solution:} $x^3 y^2 + x^2 y^4 = C$.

\vspace{0.5cm}
\textbf{20.} $(2xy - y^2) dx + (3x^2 - 4xy) dy = 0$ \\
$M' = 2x^{n+1}y^{m+1} - x^n y^{m+2} \implies M'_y = 2(m+1)x^{n+1}y^m - (m+2)x^n y^{m+1}$ \\
$N' = 3x^{n+2}y^m - 4x^{n+1}y^{m+1} \implies N'_x = 3(n+2)x^{n+1}y^m - 4(n+1)x^n y^{m+1}$ \\
Equate coefficients: \\
$2m+2 = 3n+6 \implies 3n - 2m = -4$ \\
$-(m+2) = -4n-4 \implies 4n - m = -2 \implies m = 4n+2$. \\
Substitute $m$: $3n - 2(4n+2) = -4 \implies 3n - 8n - 4 = -4 \implies -5n = 0 \implies n = 0, m = 2$. \\
Integrating factor is $\mu = y^2$. Multiply original by $y^2$: \\
$(2xy^3 - y^4) dx + (3x^2 y^2 - 4xy^3) dy = 0$. \\
\textbf{Solution:} $x^2 y^3 - x y^4 = C$.

\vspace{0.5cm}
\textbf{21.} $(2x^3 y - y^4) dx + (2xy^3 - x^4) dy = 0$ \\
$M' = 2x^{n+3}y^{m+1} - x^n y^{m+4} \implies M'_y = 2(m+1)x^{n+3}y^m - (m+4)x^n y^{m+3}$ \\
$N' = 2x^{n+1}y^{m+3} - x^{n+4}y^m \implies N'_x = 2(n+1)x^n y^{m+3} - (n+4)x^{n+3}y^m$ \\
Equating corresponding powers of $x,y$: \\
$x^{n+3}y^m$: $2m+2 = -n-4 \implies n + 2m = -6$ \\
$x^n y^{m+3}$: $-(m+4) = 2n+2 \implies 2n + m = -6$ \\
Solving system gives $m = -2$, $n = -2$. \\
$\mu = x^{-2}y^{-2}$. Multiply original by $\mu$: \\
$(2x - x^{-2}y^2) dx + (2x^{-1}y - x^2 y^{-2}) dy = 0$. \\
\textbf{Solution:} $x^2 + x^{-1}y^2 + x^2 y^{-1} = C$.

\vspace{0.5cm}
\textbf{22.} $(2y^2 + 3xy) dx + (4xy + 3x^2) dy = 0$ \\
$M' = 2x^n y^{m+2} + 3x^{n+1}y^{m+1} \implies M'_y = 2(m+2)x^n y^{m+1} + 3(m+1)x^{n+1}y^m$ \\
$N' = 4x^{n+1}y^{m+1} + 3x^{n+2}y^m \implies N'_x = 4(n+1)x^n y^{m+1} + 3(n+2)x^{n+1}y^m$ \\
Equate: $2m+4 = 4n+4 \implies m = 2n$. \\
$3m+3 = 3n+6 \implies m = n+1$. \\
$2n = n+1 \implies n = 1, m = 2$. \\
$\mu = xy^2$. Multiply original by $\mu$: \\
$(2xy^4 + 3x^2y^3) dx + (4x^2y^3 + 3x^3y^2) dy = 0$. \\
\textbf{Solution:} $x^2 y^4 + x^3 y^3 = C$.

\subsection*{C6: Standard First-Order Linear ODEs}

\textbf{23.} $\frac{dy}{dx} - 3y = e^{3x}$ \\
$P(x) = -3 \implies \mu(x) = e^{\int -3 dx} = e^{-3x}$. \\
Multiply by $e^{-3x}$: $e^{-3x} y' - 3e^{-3x} y = 1 \implies \frac{d}{dx}[e^{-3x} y] = 1$. \\
Integrate: $e^{-3x} y = x + C$. \\
\textbf{Solution:} $y = (x + C)e^{3x}$.

\vspace{0.5cm}
\textbf{24.} $x \frac{dy}{dx} + 2y = \sin(x)$ \\
Divide by $x$: $y' + \frac{2}{x}y = \frac{\sin x}{x}$. \\
$P(x) = \frac{2}{x} \implies \mu(x) = e^{\int \frac{2}{x} dx} = x^2$. \\
Multiply by $x^2$: $x^2 y' + 2xy = x\sin x \implies \frac{d}{dx}[x^2 y] = x\sin x$. \\
Integrate by parts ($u=x, dv=\sin x dx$): $\int x\sin x dx = -x\cos x + \sin x + C$. \\
\textbf{Solution:} $y = \frac{\sin x - x\cos x + C}{x^2}$.

\vspace{0.5cm}
\textbf{25.} $\frac{dy}{dx} + y \cot x = \sec x$ \\
$P(x) = \cot x \implies \mu(x) = e^{\int \cot x dx} = e^{\ln|\sin x|} = \sin x$. \\
Multiply by $\sin x$: $\sin x \, y' + \cos x \, y = \sec x \sin x \implies \frac{d}{dx}[y \sin x] = \tan x$. \\
Integrate: $y \sin x = \ln|\sec x| + C$. \\
\textbf{Solution:} $y = \csc x \ln|\sec x| + C \csc x$.

\subsection*{C7: Bernoulli Differential Equations}

\textbf{26.} $\frac{dy}{dx} + \frac{1}{x}y = x y^2$ \\
Bernoulli with $n=2$. Let $v = y^{1-2} = y^{-1}$. Then $v' = -y^{-2}y'$. \\
Multiply original by $-y^{-2}$: $-y^{-2}y' - \frac{1}{x}y^{-1} = -x \implies v' - \frac{1}{x}v = -x$. \\
Integrating factor: $\mu = e^{\int -1/x dx} = x^{-1}$. \\
Multiply by $x^{-1}$: $x^{-1}v' - x^{-2}v = -1 \implies (x^{-1}v)' = -1$. \\
$x^{-1}v = -x + C \implies v = Cx - x^2$. \\
Substitute back $v = y^{-1}$: \\
\textbf{Solution:} $y = \frac{1}{Cx - x^2}$.

\vspace{0.5cm}
\textbf{27.} $x \frac{dy}{dx} + y = y^{-2}$ \\
Divide by $x$: $y' + \frac{1}{x}y = \frac{1}{x}y^{-2}$. Bernoulli with $n = -2$. Let $v = y^3$. $v' = 3y^2y'$. \\
Multiply by $3y^2$: $3y^2y' + \frac{3}{x}y^3 = \frac{3}{x} \implies v' + \frac{3}{x}v = \frac{3}{x}$. \\
Integrating factor: $\mu = e^{\int 3/x dx} = x^3$. \\
Multiply by $x^3$: $(x^3 v)' = 3x^2 \implies x^3 v = x^3 + C \implies v = 1 + Cx^{-3}$. \\
\textbf{Solution:} $y = (1 + Cx^{-3})^{1/3}$.

\vspace{0.5cm}
\textbf{28.} $\frac{dy}{dx} - y = e^x y^2$ \\
Bernoulli with $n=2$. $v = y^{-1}$, $v' = -y^{-2}y'$. \\
Multiply by $-y^{-2}$: $-y^{-2}y' + y^{-1} = -e^x \implies v' + v = -e^x$. \\
Integrating factor: $\mu = e^x$. \\
Multiply by $e^x$: $e^x v' + e^x v = -e^{2x} \implies (e^x v)' = -e^{2x}$. \\
$e^x v = -\frac{1}{2}e^{2x} + C \implies v = Ce^{-x} - \frac{1}{2}e^x$. \\
\textbf{Solution:} $y = \frac{1}{Ce^{-x} - \frac{1}{2}e^x}$.

\subsection*{C8: Identifying Lost or Gained Solutions}

\textbf{29.} $(y^3) dx + (xy^2) dy = 0$ \\
$M_y = 3y^2, N_x = y^2$. $\frac{N_x - M_y}{M} = \frac{-2y^2}{y^3} = \frac{-2}{y} \implies \mu(y) = y^{-2}$. \\
Multiply by $y^{-2}$: $(y) dx + (x) dy = 0 \implies d(xy) = 0$. \\
Implicit Solution: $xy = C$. \\
\textbf{Analysis for lost solutions:} We multiplied the original ODE by $y^{-2}$, an operation that is undefined at $y = 0$. In the original equation, if $y = 0$, then $0 dx + 0 dy = 0$, meaning the trivial solution $y = 0$ is a valid solution. In our implicit solution $xy = C$, the line $y = 0$ cannot be generated for an arbitrary non-zero constant $C$, and algebraic division strictly forbids it. Therefore, \textbf{the solution $y = 0$ was lost}.

\subsection*{C9: Diagnostic Testing (``Find the Flaw'')}

\textbf{30.} 
\textbf{The Flaw:} The student identified $P(x)$ before putting the linear ODE into standard form (which requires the coefficient of $\frac{dy}{dx}$ to be exactly 1). \\
\textbf{Correct Step:} First divide the entire equation by $x$ to obtain $\frac{dy}{dx} + \frac{2}{x}y = x$, which establishes the correct $P(x) = \frac{2}{x}$ and leads to the integrating factor $\mu(x) = e^{\int \frac{2}{x} dx} = x^2$.

\newpage
% ---------------------------------------------------------
\section*{Appendix: Concept Checklist Reference}
% ---------------------------------------------------------
\begin{itemize}
    \item \textbf{C1: Verification of Exactness} (Questions 1, 2, 3, 4)
    \item \textbf{C2: Solving Standard Exact Equations} (Questions 5, 6, 7, 8)
    \item \textbf{C3: Finding 1D Integrating Factors ($\mu(x)$ or $\mu(y)$)} (Questions 9, 10, 11, 12)
    \item \textbf{C4: Solving Non-Exact Equations via 1D Integrating Factor} (Questions 13, 14, 15, 16, 17, 18)
    \item \textbf{C5: Algebraic Integrating Factors of form $\mu = x^n y^m$} (Questions 19, 20, 21, 22)
    \item \textbf{C6: Standard First-Order Linear ODEs} (Questions 23, 24, 25)
    \item \textbf{C7: Bernoulli Differential Equations} (Questions 26, 27, 28)
    \item \textbf{C8: Identifying Lost or Gained Solutions} (Questions 29)
    \item \textbf{C9: Diagnostic Testing (``Find the Flaw'')} (Questions 30)
\end{itemize}

\end{document}